% TOP_PROBLEM: {"id": 3346, "title": "Chowla's cosine problem", "category_id": 1, "category": {"id": 1, "name": "number_theory", "display_name": "Number Theory", "description": "Properties of integers, prime numbers, Diophantine equations.", "slug": "number-theory", "order_index": 1, "created_at": "2026-07-31T15:26:25.670Z"}, "status": "open", "problem_number": "OPG-739", "set_id": 12}
% FIRST_POSTED: 2026-10-01
% ENTRY_KIND: statement_only
% UnsolvedMath ID 3346; source catalogue code OPG-739
% !TeX program = lualatex
% Definition and sources only; this file is not an LLM solution attempt.
\documentclass[11pt,a4paper]{article}
\usepackage[margin=25mm]{geometry}
\usepackage{fontspec}
\setmainfont{lmroman10-regular.otf}[BoldFont=lmroman10-bold.otf,
 ItalicFont=lmroman10-italic.otf,BoldItalicFont=lmroman10-bolditalic.otf]
\usepackage{amsmath,amssymb,mathtools}
\usepackage{unicode-math}
\setmathfont{latinmodern-math.otf}
\AtBeginDocument{\let\setminus\smallsetminus}
\usepackage{xurl,hyperref}
\hypersetup{colorlinks=true,urlcolor=blue}
\setcounter{secnumdepth}{0}
\setlength{\parindent}{0pt}
\setlength{\parskip}{5pt}
\setlength{\emergencystretch}{3em}
\begin{document}
\section*{Chowla's cosine problem}\label{3346}
{\small UnsolvedMath ID 3346 · OPG-739 · Number Theory · OpenGarden}
\subsection{Definitions and mathematical statement}
For a nonempty finite set $A\subset\{1,2,\ldots\}$ of distinct positive
integer frequencies, define the real trigonometric polynomial
\[
 C_A(x)=\sum_{a\in A}\cos(ax),\qquad
 m(A)=-\min_{x\in[0,2\pi]}C_A(x).
\]
Continuity and periodicity ensure that this inner minimum exists. For
$n\ge1$, put
\[
 m(n)=\inf\{m(A):A\subset\mathbb N,\ |A|=n\}.
\]
The source writes an outer minimum; the infimum convention defines the
extremal value without presupposing attainment. The problem is to
determine this value, or its sharp order of growth as $n\to\infty$,
and to determine attainment if one uses the source's minimum notation.
An order-of-growth answer means matching upper and lower bounds with
absolute positive multiplicative constants, uniform in $n$.

Every coefficient is exactly one; repetitions of frequencies and a
constant term are excluded. There is no upper bound on the elements of
$A$. Since $C_A$ has average zero and $C_A(0)=n>0$, its minimum is
negative. Thus $m(A)$ measures the unavoidable negative excursion, and
the optimization asks how small that excursion can be made by choosing
the integer frequencies.
\subsection{Short English statement}
How small can the largest negative excursion of a sum of n distinct integer-frequency cosines be?
\subsection{Sources}
\begin{itemize}
\item[S1] ULAM.AI and the UnsolvedMath Contributors, supplied problems.json, record 3346 (OPG-739), OpenGarden collection. Catalogue identity and source transcription; CC BY 4.0.
\url{https://huggingface.co/datasets/ulamai/UnsolvedMath}.
\item[S2] Open Problem Garden, Chowla's cosine problem. Original problem and discussion; the mathematical formulation below is expanded from the supplied source record.
\url{http://www.openproblemgarden.org/op/chowlas_cosine_problem}.
\end{itemize}
\end{document}
